% UBT-AI-PROVENANCE-BEGIN
% schema: ubt-ai-provenance/v1
% tier: C_working
% ai_assistance: disclosed
% human_review: risk-based
% editorial_responsibility: Ing. David Jaroš
% policy: ../../../AI_PROVENANCE.md
% notice: Working material; exhaustive human review is not claimed.
% UBT-AI-PROVENANCE-END

\chapter{Program konstanty jemné struktury: co je odvozeno a co není}
\label{ch:alpha}

\section{Stav před algebrou}
Směrodatný stav vedený v repozitáři má záměrně větší váhu než jakýkoli elegantní
vzorec v této kapitole:
\begin{quote}
Fyzikální konstanta jemné struktury ještě není odvozena z prvních principů
v UBT. Cesta k celému číslu $137$ je strukturální a podmíněná; propojení, které
by určilo požadovaný efektivní koeficient z akce UBT, zůstává otevřené.
\end{quote}
Aktuálními zdroji stavu jsou \texttt{STATUS.md} a
\texttt{WHAT\_IS\_PROVED.md}. Samostatná evidence programu alfa je\\
\texttt{canonical/alpha/}\\
\texttt{ALPHA\_MASTER\_STATUS.md}. Historické texty používající
označení „odvození alfa bez přizpůsobování parametrů“ tyto evidence nepřepisují.

\section{Podmíněný potenciál vinutí}
Ústřední redukovaný model studuje potenciál
\[
 \boxed{V_{\rm eff}(n)=n^2-Bn\ln n,\qquad n>0.}
\]
Považujeme-li $n$ dočasně za spojitou proměnnou, dostaneme
\[
 \frac{dV_{\rm eff}}{dn}
 =2n-B(\ln n+1).
\]
Stacionární bod tedy splňuje
\[
 \boxed{B=\frac{2n}{\ln n+1}.}
\]
Požadujeme-li stacionární bod v $n=137$, příslušný koeficient je
\[
 B_{137}=\frac{274}{\ln137+1}\approx46.284.
\]
Tato rovnice je přesnou vlastností redukovaného potenciálu. \emph{Nevysvětluje}
však, proč musí akce UBT generovat právě tuto hodnotu $B$.

Druhá derivace je
\[
 \frac{d^2V_{\rm eff}}{dn^2}=2-\frac{B}{n}.
\]
Ve stacionárním bodě přejde do tvaru
\[
 2-\frac{2}{\ln n+1}
 =\frac{2\ln n}{\ln n+1}>0
 \qquad(n>1),
\]
takže spojitý stacionární bod je lokálním minimem. Omezení na celá čísla nebo
na prvočíselný sektor lze poté studovat jako další diskrétní výběrový problém.

\section{Co může a nemůže znamenat „prvočíselná stabilita“}
Repozitář obsahuje strukturální analýzu prvočíselné stability, v níž se
celočíselný sektor v okolí $137$ zkoumá podle stanoveného výběrového pravidla.
Minimum označené prvočíslem může být reprodukovatelným výsledkem tohoto
redukovaného modelu. Platí však:
\begin{itemize}
 \item skutečnost, že $137$ je prvočíslo, patří do standardní teorie čísel;
 \item výběr prvočíselného sektoru sám o sobě neztotožňuje
       $\alpha^{-1}$ s tímto sektorem;
 \item fyzikální korekce z $137$ na $137.036\ldots$ vyžaduje samostatné
       odvození;
 \item koeficient $B$ je stále nutné získat z dynamiky UBT, nikoli jej zpětně
       určit z cílové hodnoty.
\end{itemize}
Kapitola o větvi theta/Mellin/Feynman rozvíjí možné mechanismy, jimiž by mohla
vstoupit multiplikativní prvočíselná struktura. Tyto mechanismy však představují
výzkumná propojení, nikoli dokončené důkazy pro konstantu alfa.

\section{Audit \(N_{\rm eff}\): dvě různá počítání se nesmějí zaměňovat}
Dřívější odvození používalo
\[
 N_{\rm eff}=3\times2\times2=12.
\]
Současný audit rozlišuje alespoň dva objekty:
\begin{align*}
 N_{\rm eff}^{\rm twist}&=12,
 &&(\mathrm T),\\
 N_{\rm eff}^{\rm loop}&=3,
 &&(\mathrm L).
\end{align*}
Symbol $\mathrm T$ označuje konstrukci se zkroucením SU(2) a $\mathrm L$ přímý
auditovaný smyčkový počet skalární akce.
Identifikace
\[
 N_{\rm eff}^{\rm loop}=N_{\rm eff}^{\rm twist}
\]
není v současnosti dokázanou větou. Vzorec, který dosazuje $12$ do smyčkového
koeficientu, proto musí výslovně uvést předpoklady o zkroucení; nelze jej
popisovat jako jednoznačný důsledek samotné algebry
$\mathbb C\otimes\mathbb H$.

Repozitář proto pro dvě různé cesty eviduje dva výchozí koeficienty:
\[
 B_0^{\rm loop}=2\pi,
 \qquad
 B_0^{\rm twist}=8\pi.
\]
Přechod od kteréhokoli z nich k fenomenologicky požadované hodnotě
$B\approx46.284$ zůstává otevřený.

\section{Mezera G137-B}
Rozhodující problém programu alfa lze formulovat bez rétoriky:
\begin{quote}
Z kanonické akce UBT a jejího odůvodněného obsahu polí odvoďte efektivní
redukovaný koeficient násobící $n\ln n$, včetně všech normalizací a násobností,
a ukažte, že odvození nezávisí na naměřené hodnotě $\alpha$.
\end{quote}
Toto je mezera G137-B. Několik historických a výzkumných cest přineslo
zajímavé modulární struktury, eta a gama funkce, determinanty i struktury
vinutí. Směrodatný stav však stále vede jejich požadované promítnutí do
fyzikálního koeficientu jako otevřené.

\section{Proč má stále smysl studovat theta/Mellinovu větev}
Selhání naivního řetězce
\[
 \Theta\to\mathcal M\to\zeta\to\mathcal E\to\mathcal P
\]
kde $\mathcal M$, $\mathcal E$ a $\mathcal P$ označují Mellinovu transformaci,
Eulerův součin a prvočíselné sektory, při \emph{dynamickém} výběru prvočísel je
samo o sobě užitečné.
Eulerův součin standardní Mellinovy transformace theta funkce odráží
multiplikativní aritmetiku, která je již obsažena v množině celočíselných
indexů. Silnější mechanismus by musel dynamicky určit primitivní sektory.

Jedním z kandidátů jsou návraty theta funkce v racionálních časech, protože
kvadratické Gaussovy součty se faktorizují přes nesoudělné jmenovatele. Mocniny
prvočísel jsou v tomto popisu skutečně nerozložitelnými stavebními prvky
aritmetiky návratů. Zda lze konstrukci prvočíselného sektoru alfa přepsat v
těchto proměnných, je výslovně otevřenou otázkou, nikoli ustaveným výsledkem.

\section{Bezpečné shrnutí}
Výroky k zapamatování jsou:
\begin{enumerate}
 \item redukovaná stacionární rovnice a její výpočet stability jsou přesné;
 \item strukturální výběr celého čísla nebo prvočísla lze studovat
       reprodukovatelně;
 \item požadovaná hodnota $B$ dosud nebyla odvozena z akce UBT;
 \item evidence příslušných násobností dosud není plně sjednocena;
 \item $\alpha^{-1}=137.036\ldots$ dnes není výsledkem UBT odvozeným z prvních
       principů.
\end{enumerate}
Toto shrnutí je záměrně méně dramatické než starší archivní text a pro budoucí
práci užitečnější: přesně ukazuje, kterou rovnici je nutné uzavřít jako další.
